% Part: sets-functions-relations
% Chapter: functions
% Section: composition

\documentclass[../../../include/open-logic-section]{subfiles}

\begin{document}

\olfileid{sfr}{fun}{cmp}
\olsection{فنکشناں دی ترکیب (Composition of Functions)}

\begin{explain}
\oliflabeldef{sfr:fun:inv:sec}{اسی \olref[inv]{sec} وچ ویکھیا سی
کہ اُلٹ~$f^{-1}$، !!a{bijection} (بائجیکشن؛ bijection)~$f$ دا، آپ وی
اک فنکشن اے۔ فنکشناں اُتے اک ہور عمل فنکشن دی ترکیب (function
composition) اے: ا}{ا}سی دو فنکشناں، $f$ تے~$g$، دی ترکیب نال نواں
فنکشن تعریف کر سکدے آں: پہلے~$f$ تے اوہدے بعد~$g$ لاگو کریے۔ ایہ
اوہدوں ای ممکن اے جدوں حاصل قدراں دے سیٹ (ranges) تے دائرۂ تعریف
(domains) آپس وچ میل کھان؛ یعنی~$f$ دا حاصل قدراں دا سیٹ (range)،
~$g$ دے دائرۂ تعریف (domain) دا ذیلی سیٹ (subset) ہونا چاہیدا اے۔
\oliflabeldef{sfr:rel:ops:sec}{فنکشناں اُتے ایہ عمل، تعلقاں (relations)
اُتے \olref[rel][ops]{sec} وچ دتے نسبتی حاصل ضرب (relative product)
دے عمل دا مماثل (analogue) اے۔}{}

اک خاکہ فنکشناں دی ترکیب (composition) دا خیال سمجھاؤن وچ مدد دے
سکدا اے۔ \olref{fig:composition} وچ اسی دو فنکشن $f \colon A \to B$
تے $g \colon B \to C$، تے اوہناں دی ترکیب~$(\comp{f}{g})$ وکھاؤندے
آں۔ فنکشن $(\comp{f}{g}) \colon A \to C$،~$A$ دے ہر !!{element}
(element/member) نوں~$C$ دے !!a{element} نال جوڑدا اے۔~$C$ دے کیہڑے
!!{element} نال~$A$ دا !!a{element} جوڑیا جاندا اے، ایہ ایویں
طے ہوندا اے: انپٹ (input) $x \in A$ دتا ہووے، تاں پہلے فنکشن~$f$
نوں~$x$ اُتے لاگو کرو؛ اوہ $f(x) = y \in B$ آؤٹ پٹ (output) دیندا
اے۔ فیر فنکشن~$g$ نوں~$y$ اُتے لاگو کرو؛ اوہ
$g(f(x)) = g(y) = z \in C$ آؤٹ پٹ دیندا اے۔
\begin{figure}
  \olasset[2\olphotowidth]{assets/diagrams/composition.tikz}
  \caption{ترکیب (composition) $g \circ f$، دو فنکشناں $f$ تے~$g$ دی۔}
  \ollabel{fig:composition}
\end{figure}
\end{explain}

\begin{defn}[ترکیب (Composition)]
من لو کہ $f\colon A \to B$ تے $g\colon B \to C$ فنکشن نیں۔ پہلے
$f$ تے فیر~$g$ لا کے بنائی \emph{ترکیب} (composition)،
$\comp{f}{g} \colon A \to C$ اے، جتھے
$(\comp{f}{g})(x) = g(f(x))$۔
\end{defn}

\begin{ex}
فنکشن $f(x) = x + 1$ تے $g(x) = 2x$ اُتے غور کرو۔ کیوں جے
$(\comp{f}{g})(x) = g(f(x))$، ہر انپٹ (input)~$x$ لئی پہلے اوہدا
جانشین (successor) لینا، تے فیر حاصل نوں دو نال ضرب دینا اے۔ ایس لئی
اوہناں دی ترکیب (composition) $(\comp{f}{g})(x) = 2(x+1)$ نال دتی
جاندی اے۔
\end{ex}

\begin{prob}
وکھاؤ کہ جے $f \colon A \to B$ تے $g \colon B \to C$ دونیں
!!{injective} (اک بہ اک؛ injective) ہون، تاں
$\comp{f}{g}\colon A \to C$ وی !!{injective} (اک بہ اک؛ injective)
اے۔
\end{prob}

\begin{prob}
وکھاؤ کہ جے $f \colon A \to B$ تے $g \colon B \to C$ دونیں
!!{surjective} (ہر ہدف قدر لین والے؛ surjective) ہون، تاں
$\comp{f}{g}\colon A \to C$ وی !!{surjective} (ہر ہدف قدر لین والا؛
surjective) اے۔
\end{prob}

\begin{prob}
من لو کہ $f \colon A \to B$ تے $g \colon B \to C$۔ وکھاؤ کہ
ترکیبی فنکشن~$\comp{f}{g}$ دا گراف (graph) $R_f \mid R_g$ اے۔
\end{prob}

\end{document}
